\lecture{11}{Sep 16 10:15}{Gauss-Jordan Elimination}
\section{Gauss-Jordan Elimination}
\begin{definition}[RREF]\label{def:rref}
    A matrix is in \textbf{reduced row echelon form} (RREF) if it is in REF and:
    \begin{enumerate}
        \item All leading coefficients are $\tilde{1}$
        \item In a column that includes a leading coefficient, all other coefficients are $\tilde{0}$
    \end{enumerate}
\end{definition}
\begin{note}
    This definition applies in every field
\end{note}
Given a matrix of order $m \times n$
\begin{enumerate}
    \item Consider row $i=1$
    \item Identify the first (leftmost) column containing a non-zero entry on or below row $i$, column $j$
    \item If entry $a_{ij}=0$, swap row $i$ with any row $k>i$ that doesn't have a zero in column $j$
    \item Now that $a_{ij}\neq 0$, multiply row $i$ by $a_{ij}^{-1}$
    \item If there are any non-zero entries left in column $j$, add multiples of row $i$ to the rows of those entries to make them zero.
    \item Increase $i \to i+1$ and repeat steps $2-5$
\end{enumerate}

However, for solving a system of linear equations, RREF is not worth it (REF is faster and easy enough ($\frac{1}{2}O(n^3)$ vs $\frac{1}{3}O(n^3)+\frac{1}{2}O(n^2)$))